\sectionkicker{Source-backed technical notes}
\subsection*{Runtimeは同じ月に何を吸収したか — Transformers v5・vLLM・SGLang: Technical Notes}
本欄は記事本文で圧縮した一次資料上の情報を、比較・再検証しやすい形で整理したものである。時系列、確認済みの主張、留意点、一次資料URLを掲載する。
\medskip
\subsection*{Theme at a glance}
\addcontentsline{toc}{subsection}{Theme at a glance}
\begin{center}
\footnotesize
\begin{tabularx}{\linewidth}{@{}p{0.20\linewidth}p{0.10\linewidth}p{0.14\linewidth}X@{}}
\toprule
資料 & 位置づけ & 種別 & 時系列 \\
\midrule
Transformers v5 and v5.0.0rc3 & 主要資料 & Framework & 2026-01-26 (プロジェクトPre-release); 2026-01-26 (プロジェクト公開) \\
vLLM v0.14.0 and v0.15.0 & 主要資料 & Framework & 2026-01-20 (プロジェクト公開); 2026-01-29 (プロジェクト公開) \\
SGLang v0.5.8 & 補足資料 & Framework & 2026-01-23 (プロジェクト公開) \\
\bottomrule
\end{tabularx}
\end{center}
\normalsize

\begin{technicalnote}{Transformers v5 and v5.0.0rc3}{主要資料}
\begingroup
% reader-facing Technical Notes break policy
\widowpenalty=10000
\clubpenalty=10000
\displaywidowpenalty=10000
\begin{tabularx}{\linewidth}{@{}>{\bfseries}p{0.22\linewidth}X@{}}
組織 & Hugging Face \\
種別 & Framework \\
時系列 & 2026-01-26 (プロジェクトPre-release); 2026-01-26 (プロジェクト公開) \\
\end{tabularx}
\smallskip
{\bfseries 一次資料から整理したtechnical points}
\begin{itemize}[leftmargin=1.5em,itemsep=0.35em]
\item \textbf{一次情報で確認できる事実}: Hugging Faceは1月26日、5年ぶりのmajor releaseとなるTransformers v5を公開し、新しいdynamic weight-loading／WeightConverter設計やtokenizer変更など大きなAPI変更を導入した。
\item \textbf{一次情報で確認できる事実}: v5 release notesはbreaking changeとquantization／loadingの整理、deprecated loading pathからのmigration、広範なarchitecture refactoringを説明する。
\item \textbf{一次情報で確認できる事実}: v5.0.0rc3 notesはstable release直前にGLM-4.7／GLM-ImageとMiniMax-M2 supportを追加し、新しいmodel familyが迅速に統合される様子を示す。
\end{itemize}
\begin{minipage}{\linewidth}
% reader-facing Technical Notes generic boundary/source tail group
{\bfseries 読む際の境界}
\begin{itemize}[leftmargin=1.5em,itemsep=0.35em]
\item \textbf{分析上の留意点}: RCとstable releaseは同日であり、stable v5を主要event、rc3をmodel integrationのsupporting evidenceとして扱う。
\item \textbf{分析上の留意点}: library supportが示すのはimplementation compatibilityであり、endorsementや独立測定されたmodel qualityではない。
\end{itemize}
\begin{samepage}
{\bfseries 一次資料}
\begingroup\sloppy
\Urlmuskip=0mu plus 2mu\relax
\begin{itemize}[leftmargin=1.5em,itemsep=0.25em]
\item {\scriptsize\url{https://github.com/huggingface/transformers/releases/tag/v5.0.0}}
\item {\scriptsize\url{https://github.com/huggingface/transformers/releases/tag/v5.0.0rc3}}
\end{itemize}
\endgroup
\end{samepage}
\end{minipage}
\endgroup
\end{technicalnote}
\medskip

\begin{technicalnote}{vLLM v0.14.0 and v0.15.0}{主要資料}
\begingroup
% reader-facing Technical Notes break policy
\widowpenalty=10000
\clubpenalty=10000
\displaywidowpenalty=10000
\begin{tabularx}{\linewidth}{@{}>{\bfseries}p{0.22\linewidth}X@{}}
組織 & vLLM Project \\
種別 & Framework \\
時系列 & 2026-01-20 (プロジェクト公開); 2026-01-29 (プロジェクト公開) \\
\end{tabularx}
\smallskip
{\bfseries 一次資料から整理したtechnical points}
\begin{itemize}[leftmargin=1.5em,itemsep=0.35em]
\item \textbf{一次情報で確認できる事実}: vLLM v0.14.0はasync schedulingをdefaultで有効化し、engine、tool-calling、API、model supportを多数更新した。release notesはこれをbreaking changeとし、unsupported configurationも明記する。
\item \textbf{一次情報で確認できる事実}: vLLM v0.15.0はnew model supportにKimi-K2.5 architectureを追加し、async schedulingをpipeline parallelismへ拡張したほか、streaming input、quantization、platformも更新した。
\item \textbf{Project claim}: release notes内のspeedup値はproject報告であり、configurationに依存する。
\end{itemize}
\begin{minipage}{\linewidth}
% reader-facing Technical Notes generic boundary/source tail group
{\bfseries 読む際の境界}
\begin{itemize}[leftmargin=1.5em,itemsep=0.35em]
\item \textbf{分析上の留意点}: seriesは9日間隔の二つのreleaseを含むため、一つのatomic featureではなくserving／runtime evolutionとして要約する。
\item \textbf{分析上の留意点}: projectのperformance claimは独立再現していない。
\end{itemize}
\begin{samepage}
{\bfseries 一次資料}
\begingroup\sloppy
\Urlmuskip=0mu plus 2mu\relax
\begin{itemize}[leftmargin=1.5em,itemsep=0.25em]
\item {\scriptsize\url{https://github.com/vllm-project/vllm/releases/tag/v0.14.0}}
\item {\scriptsize\url{https://github.com/vllm-project/vllm/releases/tag/v0.15.0}}
\end{itemize}
\endgroup
\end{samepage}
\end{minipage}
\endgroup
\end{technicalnote}
\medskip

\begin{technicalnote}{SGLang v0.5.8}{補足資料}
\begingroup
% reader-facing Technical Notes break policy
\widowpenalty=10000
\clubpenalty=10000
\displaywidowpenalty=10000
\begin{tabularx}{\linewidth}{@{}>{\bfseries}p{0.22\linewidth}X@{}}
組織 & SGLang Project \\
種別 & Framework \\
時系列 & 2026-01-23 (プロジェクト公開) \\
\end{tabularx}
\smallskip
{\bfseries 一次資料から整理したtechnical points}
\begin{itemize}[leftmargin=1.5em,itemsep=0.35em]
\item \textbf{一次情報で確認できる事実}: SGLang v0.5.8はGLM-4.7-Flashのday-zero support、diffusion serving拡張、FlashAttention 4 decoding、複数のlong-context／disaggregated-serving変更を追加した。
\item \textbf{Project claim}: 最大1.5倍のdiffusion高速化、chunked pipelineのnear-linear scaling、特定GLM4-MoE構成で65\%のTTFT改善などはprojectのrelease noteによる報告値である。
\end{itemize}
\begin{minipage}{\linewidth}
% reader-facing Technical Notes generic boundary/source tail group
{\bfseries 読む際の境界}
\begin{itemize}[leftmargin=1.5em,itemsep=0.35em]
\item \textbf{分析上の留意点}: releaseは多数のupstream PRと、異なるmodel／hardware条件のperformance claimを集約している。
\item \textbf{分析上の留意点}: model supportは実装上利用可能になったことを示すが、同等の品質やproduction validationを意味しない。
\end{itemize}
\begin{samepage}
{\bfseries 一次資料}
\begingroup\sloppy
\Urlmuskip=0mu plus 2mu\relax
\begin{itemize}[leftmargin=1.5em,itemsep=0.25em]
\item {\scriptsize\url{https://github.com/sgl-project/sglang/releases/tag/v0.5.8}}
\end{itemize}
\endgroup
\end{samepage}
\end{minipage}
\endgroup
\end{technicalnote}
\medskip
